% !Mode:: "TeX:UTF-8"
% !TEX root = ../../main.tex
% ============================================================
% TOPIC 23
% ============================================================
\mashgulot{23}{The iterative method for the approximate solution of systems of linear algebraic equations}

For large systems the direct methods (Cramer’s rule, Gaussian elimination)
require a very large number of operations, and rounding errors accumulate.
In such cases one uses \textbf{iterative methods}, which produce the solution
as the limit of a sequence of successive approximations.

\subsection{The simple iteration (Jacobi) method}

\begin{tarif}
Solving the $i$-th equation of the system $Ax=b$ for $x_i$ (we assume
$a_{ii}\neq0$), we bring the system to the form $x=Bx+c$:
\[\kf{x_i=\frac{1}{a_{ii}}\Big(b_i-\sum\limits_{j\neq i}a_{ij}x_j\Big),\qquad
b_{ij}=-\frac{a_{ij}}{a_{ii}}\ (j\neq i),\quad b_{ii}=0,\quad c_i=\frac{b_i}{a_{ii}} .}\]
We choose an initial approximation $x^{(0)}$ (usually $x^{(0)}=c$ or
$x^{(0)}=0$) and compute the successive approximations by the formula
\[\kf{x^{(k+1)}=Bx^{(k)}+c,\qquad k=0,1,2,\dots}\]
\end{tarif}

\begin{teorema}[A sufficient condition for convergence]
If $\|B\|<1$ (for instance,
$\|B\|_\infty=\max\limits_i\sum\limits_j|b_{ij}|<1$), then the iteration
converges to the unique solution $x^*$ for every choice of $x^{(0)}$, and
\[\kf{\|x^{(k)}-x^*\|\leq\frac{\|B\|^k}{1-\|B\|}\,\|x^{(1)}-x^{(0)}\|.}\]
In particular, the method converges whenever $A$ is \textbf{diagonally dominant},
that is, $|a_{ii}|>\sum\limits_{j\neq i}|a_{ij}|$ in every row.
\end{teorema}

\begin{misol}
Solve the system
$\begin{cases}10x_1+x_2+x_3=12,\\2x_1+10x_2+x_3=13,\\2x_1+2x_2+10x_3=14\end{cases}$
by the simple iteration method.

\textbf{Solution.} The matrix is diagonally dominant ($10>2$, $10>3$, $10>4$).
Solving each equation for its diagonal unknown, we obtain
\[\kf{\begin{cases}x_1=1.2-0.1x_2-0.1x_3,\\
x_2=1.3-0.2x_1-0.1x_3,\\
x_3=1.4-0.2x_1-0.2x_2,\end{cases}\qquad
\|B\|_\infty=\max(0.2;\,0.3;\,0.4)=0.4<1 .}\]
Starting from $x^{(0)}=(0;0;0)$, we get the following approximations:

\begin{center}
\begin{tabular}{c|ccc}
$k$ & $x_1^{(k)}$ & $x_2^{(k)}$ & $x_3^{(k)}$\\ \hline
0 & 0 & 0 & 0\\
1 & 1.2 & 1.3 & 1.4\\
2 & 0.93 & 0.92 & 0.90\\
3 & 1.018 & 1.024 & 1.030\\
\end{tabular}
\end{center}
The approximations oscillate about the exact solution $x^*=(1;1;1)$; after
three iterations the error does not exceed $0.03$.
\end{misol}

\begin{eslatma}
If the matrix is not diagonally dominant, it is advisable to bring the system
to a diagonally dominant form by interchanging the equations or by replacing
them with suitable linear combinations.
\end{eslatma}

\subsection{Iterative methods in information security}

\textbf{Large sparse systems.} In cryptanalysis (for example, in the quadratic
sieve and in index calculus methods) one meets sparse systems with hundreds of
thousands of unknowns. Gaussian elimination is out of the question there, both
in memory and in time, so iterative methods (Lanczos, Wiedemann and others)
are used instead.

\textbf{Reputation and PageRank.} The trust scores of the nodes of a network
form the solution of the equation $r=dPr+(1-d)v$. It is found by the simple
iteration $r^{(k+1)}=dPr^{(k)}+(1-d)v$, which converges because $d<1$.

\begin{misol}
Let the trust scores in a network with three nodes satisfy
$r_1=0.5r_2+0.2$, \ $r_2=0.5r_3+0.2$, \ $r_3=0.5r_1+0.2$.
Starting from $r^{(0)}=(0;0;0)$, we get $r^{(1)}=(0.2;0.2;0.2)$,
$r^{(2)}=(0.3;0.3;0.3)$, $r^{(3)}=(0.35;0.35;0.35)$, \dots, while the exact
solution is $r_i=0.4$.
\end{misol}

\subsection{Computing with Python, Maple and Mathcad}

The iteration of the example above takes only a few lines of NumPy. We form
$B$ and $c$ from $A$ and $b$, check the sufficient condition and perform
three iterations:
\begin{Verbatim}[fontsize=\footnotesize,frame=single,label={\small Python}]
import numpy as np
A = np.array([[10, 1, 1], [2, 10, 1], [2, 2, 10]],
             dtype=float)
b = np.array([12, 13, 14], dtype=float)
d = np.diag(A)
B = -A / d[:, None]
np.fill_diagonal(B, 0)
c = b / d
print("||B||_inf =", np.abs(B).sum(axis=1).max())
x = np.zeros(3)
for k in range(1, 4):
    x = B @ x + c
    print(k, x.round(4))
\end{Verbatim}
\begin{Verbatim}[fontsize=\footnotesize,frame=single,label={\small Output}]
||B||_inf = 0.4
1 [1.2 1.3 1.4]
2 [0.93 0.92 0.9 ]
3 [1.018 1.024 1.03 ]
\end{Verbatim}
In Maple we write $B=E-D^{-1}A$ and $c=D^{-1}b$, where $D$ is the diagonal
part of $A$ (the name \texttt{D} is reserved in Maple, so we use \texttt{D1}):
\begin{Verbatim}[fontsize=\footnotesize,frame=single,label={\small Maple}]
with(LinearAlgebra):
A := Matrix([[10, 1, 1], [2, 10, 1], [2, 2, 10]]):
b := Vector([12, 13, 14]):
D1 := DiagonalMatrix(Diagonal(A)):
B := IdentityMatrix(3) - D1^(-1) . A:
c := D1^(-1) . b:
Norm(B, infinity);             # 2/5
x := Vector(3):
for k to 3 do x := B . x + c end do:
evalf(x);                      # 1.018, 1.024, 1.030
\end{Verbatim}
In Mathcad (with \texttt{ORIGIN} $=0$) the iteration can be written with a
range variable: each column of the matrix $X$ is one approximation.
The function \texttt{lsolve} gives the exact solution for comparison, and
\texttt{normi} is the $\infty$-norm:
\begin{Verbatim}[fontsize=\footnotesize,frame=single,label={\small Mathcad}]
     | 10  1  1 |        | 12 |
A := |  2 10  1 |   b := | 13 |
     |  2  2 10 |        | 14 |
i := 0..2     d_i := A_(i,i)     D := diag(d)
B := identity(3) - D^-1 * A      c := D^-1 * b
normi(B) = 0.4
X<0> := 0 * c     k := 0..2
X<k+1> := B * X<k> + c
          | 0   1.2   0.93   1.018 |
     X =  | 0   1.3   0.92   1.024 |
          | 0   1.4   0.90   1.030 |
lsolve(A, b) = (1 1 1)^T
\end{Verbatim}

\subsection{Exercises on the simple iteration method}

\emph{1. Perform two iterations of the simple iteration (Jacobi) method for the system, starting from $x^{(0)}=(0;0;0)$, and compare the result with the exact solution.}

\bigskip
\begin{enumerate}[label=\arabic*., leftmargin=*]
\item $\begin{cases}-8x_1-x_2+3x_3=15\\-3x_1+9x_2+x_3=15\\-x_1-4x_3=2\end{cases}$
\item $\begin{cases}-11x_1-3x_2-3x_3=-19\\2x_1+6x_2=-8\\-3x_1+x_2+9x_3=1\end{cases}$
\item $\begin{cases}-6x_1+2x_2-x_3=-6\\x_1-8x_2+3x_3=-20\\-x_1+x_2-8x_3=16\end{cases}$
\item $\begin{cases}8x_1+3x_2-2x_3=-1\\x_1+9x_2+2x_3=-2\\3x_1+x_2+9x_3=29\end{cases}$
\item $\begin{cases}9x_1-x_2-3x_3=-20\\-x_1+5x_2+x_3=-2\\-x_1+3x_2+10x_3=9\end{cases}$
\item $\begin{cases}-9x_1+2x_2+x_3=26\\-3x_1+9x_2-2x_3=29\\-3x_2+6x_3=3\end{cases}$
\item $\begin{cases}-5x_1+x_2=-7\\2x_1+6x_2-x_3=-13\\-3x_1+3x_2-12x_3=-45\end{cases}$
\item $\begin{cases}8x_1-x_2-2x_3=-1\\-3x_1+10x_2+x_3=29\\x_1+2x_2+6x_3=0\end{cases}$
\item $\begin{cases}8x_1+3x_3=15\\-2x_1-10x_2+2x_3=-2\\3x_1-6x_3=27\end{cases}$
\item $\begin{cases}7x_1+2x_2-2x_3=17\\-x_1-6x_2+2x_3=1\\3x_1+9x_3=27\end{cases}$
\item $\begin{cases}-8x_1-3x_2-x_3=22\\7x_2+3x_3=6\\-2x_1+2x_2-10x_3=-14\end{cases}$
\item $\begin{cases}-6x_1+x_3=-3\\2x_1+8x_2+2x_3=18\\x_1-2x_2+9x_3=-33\end{cases}$
\item $\begin{cases}-6x_1+x_3=2\\-6x_2+2x_3=-14\\-x_1-4x_3=-8\end{cases}$
\item $\begin{cases}9x_1+3x_2+x_3=-29\\-x_1+7x_2-x_3=-5\\2x_1-3x_2+11x_3=8\end{cases}$
\item $\begin{cases}9x_1-3x_3=-21\\3x_1+10x_2-x_3=-17\\3x_1+3x_2+11x_3=2\end{cases}$
\item $\begin{cases}-9x_1+x_2+3x_3=-29\\3x_1+9x_2-2x_3=-6\\2x_1+x_2+6x_3=-16\end{cases}$
\item $\begin{cases}6x_1+x_2+2x_3=-13\\-2x_1+8x_2-x_3=-20\\2x_2+6x_3=-18\end{cases}$
\item $\begin{cases}-9x_1-2x_2-2x_3=3\\8x_2+3x_3=9\\x_1-2x_2-7x_3=-22\end{cases}$
\item $\begin{cases}10x_1+3x_2-3x_3=-39\\-7x_2-x_3=5\\2x_1-2x_2-7x_3=-18\end{cases}$
\item $\begin{cases}8x_1+3x_3=-1\\-2x_1-6x_2=-2\\-x_1-2x_2-8x_3=23\end{cases}$
\item $\begin{cases}-6x_1+x_2+2x_3=7\\-4x_2+x_3=14\\-2x_1-8x_3=-14\end{cases}$
\item $\begin{cases}-4x_1=12\\-x_1+6x_2+2x_3=3\\2x_1+2x_2-7x_3=-6\end{cases}$
\item $\begin{cases}8x_1-2x_2+2x_3=2\\3x_1-11x_2-3x_3=-3\\x_1-3x_2-9x_3=-9\end{cases}$
\item $\begin{cases}5x_1+x_2-x_3=-13\\-x_1+7x_2-3x_3=21\\-2x_1+3x_2-10x_3=5\end{cases}$
\item $\begin{cases}-11x_1+3x_2-3x_3=-7\\3x_1-8x_2+2x_3=-22\\2x_2+5x_3=-4\end{cases}$
\item $\begin{cases}11x_1+3x_2-2x_3=-34\\8x_2+2x_3=-10\\2x_1-2x_2+7x_3=21\end{cases}$
\item $\begin{cases}-6x_1-x_2+2x_3=-20\\x_1-7x_2+x_3=15\\-3x_2-6x_3=18\end{cases}$
\item $\begin{cases}12x_1+3x_2+3x_3=-3\\3x_1+9x_2-2x_3=13\\-2x_1-2x_2+7x_3=5\end{cases}$
\item $\begin{cases}-9x_1+3x_2+x_3=-21\\2x_1+9x_2+3x_3=24\\3x_1-2x_2-10x_3=33\end{cases}$
\item $\begin{cases}9x_1-2x_2-2x_3=-27\\-x_1+4x_2=7\\-2x_2+6x_3=-8\end{cases}$
\end{enumerate}

\emph{2. Check whether the sufficient condition for the convergence of the simple iteration method ($\|B\|_\infty<1$) holds.}
\bigskip
\begin{multicols}{2}\raggedright
\begin{enumerate}[label=\arabic*., leftmargin=*]
\item $A=\begin{pmatrix}-10 & -3 & 3\\ -3 & -11 & 3\\ -3 & -2 & -8\end{pmatrix}$
\item $A=\begin{pmatrix}10 & 3 & -3\\ 0 & 5 & 2\\ 0 & 0 & 5\end{pmatrix}$
\item $A=\begin{pmatrix}3 & 2 & -2\\ 2 & 2 & -2\\ 4 & 0 & -2\end{pmatrix}$
\item $A=\begin{pmatrix}-6 & 2 & -1\\ 2 & -11 & 3\\ -3 & -1 & -7\end{pmatrix}$
\item $A=\begin{pmatrix}11 & 3 & -3\\ -1 & -8 & -3\\ -3 & 1 & 9\end{pmatrix}$
\item $A=\begin{pmatrix}2 & 4 & -2\\ 3 & 3 & 2\\ 1 & 1 & 1\end{pmatrix}$
\item $A=\begin{pmatrix}4 & 0 & -1\\ 1 & 6 & 2\\ -3 & -3 & -11\end{pmatrix}$
\item $A=\begin{pmatrix}-9 & -3 & 2\\ -1 & -7 & -1\\ 1 & 3 & 8\end{pmatrix}$
\item $A=\begin{pmatrix}3 & 3 & 4\\ 2 & 2 & -2\\ 3 & -4 & 2\end{pmatrix}$
\item $A=\begin{pmatrix}9 & 2 & -2\\ 2 & -6 & -1\\ 3 & 1 & -9\end{pmatrix}$
\item $A=\begin{pmatrix}10 & 3 & 1\\ 0 & 9 & 3\\ -2 & -3 & 10\end{pmatrix}$
\item $A=\begin{pmatrix}2 & -1 & -1\\ -4 & 3 & 4\\ -1 & 2 & 2\end{pmatrix}$
\item $A=\begin{pmatrix}11 & -3 & 3\\ 1 & 7 & -3\\ 2 & -3 & -10\end{pmatrix}$
\item $A=\begin{pmatrix}9 & -3 & 3\\ -2 & 8 & -3\\ 3 & -2 & 9\end{pmatrix}$
\item $A=\begin{pmatrix}2 & 4 & 0\\ 3 & 1 & 4\\ -4 & -4 & -2\end{pmatrix}$
\item $A=\begin{pmatrix}9 & 3 & 1\\ 0 & 8 & 3\\ -3 & 1 & 10\end{pmatrix}$
\item $A=\begin{pmatrix}-9 & -1 & 3\\ 2 & 9 & -1\\ 2 & 2 & 7\end{pmatrix}$
\item $A=\begin{pmatrix}2 & 2 & 2\\ 2 & 3 & -3\\ 4 & -2 & 1\end{pmatrix}$
\item $A=\begin{pmatrix}-6 & -3 & 0\\ 1 & 7 & -3\\ 0 & -1 & -5\end{pmatrix}$
\item $A=\begin{pmatrix}8 & 2 & 2\\ 1 & -9 & -2\\ -2 & -1 & 9\end{pmatrix}$
\item $A=\begin{pmatrix}3 & 2 & -1\\ -1 & 2 & -2\\ 3 & -4 & 2\end{pmatrix}$
\item $A=\begin{pmatrix}9 & 3 & 3\\ 1 & -6 & 0\\ 3 & 1 & -9\end{pmatrix}$
\item $A=\begin{pmatrix}8 & -1 & -2\\ 2 & -8 & -1\\ 3 & 3 & 10\end{pmatrix}$
\item $A=\begin{pmatrix}3 & -4 & 1\\ -3 & 1 & -3\\ 4 & -3 & -2\end{pmatrix}$
\item $A=\begin{pmatrix}8 & -1 & 1\\ 2 & 9 & -2\\ -1 & -1 & -5\end{pmatrix}$
\item $A=\begin{pmatrix}-8 & -2 & -2\\ 1 & 9 & 2\\ 0 & 1 & -7\end{pmatrix}$
\item $A=\begin{pmatrix}1 & -4 & 3\\ -3 & 3 & -3\\ -4 & 0 & 2\end{pmatrix}$
\item $A=\begin{pmatrix}-9 & 2 & -1\\ 0 & -7 & 2\\ 1 & 3 & 10\end{pmatrix}$
\item $A=\begin{pmatrix}-5 & 0 & -2\\ -2 & -10 & 2\\ 2 & 0 & 6\end{pmatrix}$
\item $A=\begin{pmatrix}3 & -3 & -1\\ -4 & -2 & 0\\ 1 & -3 & 1\end{pmatrix}$
\end{enumerate}
\end{multicols}

\emph{3. Applied problems (10 problems).}
\bigskip
\begin{enumerate}[label=\arabic*., leftmargin=*]
\item \textbf{(Three iterations.)} Perform three iterations of the simple iteration method for the system $\begin{cases}8x_1+3x_2=25\\3x_1+6x_2=24\end{cases}$, starting from $x^{(0)}=(0;0)$, and compare the result with the exact solution.
\item \textbf{(Three iterations.)} Perform three iterations of the simple iteration method for the system $\begin{cases}8x_1+3x_2=27\\-2x_1+5x_2=-1\end{cases}$, starting from $x^{(0)}=(0;0)$, and compare the result with the exact solution.
\item \textbf{(Three iterations.)} Perform three iterations of the simple iteration method for the system $\begin{cases}8x_1-3x_2=22\\-2x_1+5x_2=-14\end{cases}$, starting from $x^{(0)}=(0;0)$, and compare the result with the exact solution.
\item \textbf{(Three iterations.)} Perform three iterations of the simple iteration method for the system $\begin{cases}5x_1+x_2=-23\\3x_1+7x_2=-33\end{cases}$, starting from $x^{(0)}=(0;0)$, and compare the result with the exact solution.
\item \textbf{(Three iterations.)} Perform three iterations of the simple iteration method for the system $\begin{cases}7x_1+3x_2=19\\3x_1+5x_2=-3\end{cases}$, starting from $x^{(0)}=(0;0)$, and compare the result with the exact solution.
\item \textbf{(Three iterations.)} Perform three iterations of the simple iteration method for the system $\begin{cases}9x_1+x_2=-6\\-x_1+6x_2=19\end{cases}$, starting from $x^{(0)}=(0;0)$, and compare the result with the exact solution.
\item \textbf{(Three iterations.)} Perform three iterations of the simple iteration method for the system $\begin{cases}9x_1-x_2=-17\\-x_1+9x_2=-7\end{cases}$, starting from $x^{(0)}=(0;0)$, and compare the result with the exact solution.
\item \textbf{(Three iterations.)} Perform three iterations of the simple iteration method for the system $\begin{cases}7x_1-3x_2=-6\\-2x_1+9x_2=18\end{cases}$, starting from $x^{(0)}=(0;0)$, and compare the result with the exact solution.
\item \textbf{(Three iterations.)} Perform three iterations of the simple iteration method for the system $\begin{cases}5x_1+2x_2=18\\3x_1+8x_2=4\end{cases}$, starting from $x^{(0)}=(0;0)$, and compare the result with the exact solution.
\item \textbf{(Three iterations.)} Perform three iterations of the simple iteration method for the system $\begin{cases}6x_1+2x_2=32\\-3x_1+5x_2=8\end{cases}$, starting from $x^{(0)}=(0;0)$, and compare the result with the exact solution.
\end{enumerate}

